\documentclass[../main.tex]{subfiles}

\begin{document}

\begin{motivation}
	\textbf{「附录是武器库，不是新战场」}——当数学形式扑面而来时，我们容易迷失在符号丛林，忘记它只是工具而非目的。
	
	\textbf{核心动机}：防止读者误将数学附录视为新的叙事主线，明确其「形式化重写」的定位，让附录成为加固理解的脚手架，而非分散注意力的装饰。
\end{motivation}
	
	\begin{logicmap}
	\begin{tikzpicture}[node distance=1.2cm, every node/.style={draw, rectangle, rounded corners, align=center}]
	\node (problem) {数学符号丛林：迷失在工具中};
	\node (insight) [below left=of problem] {附录定位：武器库而非新战场};
	\node (guide) [below right=of problem] {使用方式：按需跳转，与正文一一对应};
	\draw[-{Stealth}] (problem) -- (insight);
	\draw[-{Stealth}] (problem) -- (guide);
	\end{tikzpicture}
	\end{logicmap}
	
	\section{V.0 使用说明与工程声明}
	
	\begin{questionbox}
\paragraph*{读者疑问}
为什么一本讨论社会动力学的书需要数学附录？附录会不会变成另一个「需要学习的全新内容」，反而偏离主题？
\end{questionbox}

	% 动机说明
	% 本节目的：防止读者误将数学附录视为新的叙事主线
	\begin{mathbox}{附录定位与阅读方式}
		\textbf{定位}：本附录不是新的论述层，而是对正文关键论断的\textbf{形式化、可核验重写}。
		
		\textbf{阅读建议}：
		\begin{itemize}
			\item 可在阅读正文后，按需跳转至对应数学模块；
			\item 每一节均通过 \texttt{\mathlink} 与正文一一对应；
			\item 附录中不引入新的经验性主张，仅展开假设、模型与边界。
		\end{itemize}
	\end{mathbox}
	
	\paragraph*{逻辑衔接}
本附录通过形式化手段加固正文的论证链条，但它本身不是论证的终点。每一个数学模块都指向正文中的具体论断，读者应始终在「数学细节」与「直观解释」之间来回切换。附录的价值不在于 mastering the math，而在于 seeing the structure。
	
\end{document}
